% ===========================================================================
%  preamble.tex --- packages and macros for the population chapter.
%  When this chapter is dropped into the thesis, merge the contents of this
%  file into the thesis preamble; duplicates may simply be deleted.
% ===========================================================================

\usepackage[T1]{fontenc}
\usepackage[utf8]{inputenc}
\usepackage{lmodern}
\usepackage{microtype}
\usepackage{amsmath,amssymb,amsthm,mathtools}
\usepackage{graphicx}
\usepackage{float}
\usepackage{placeins}
\usepackage{booktabs}
\usepackage[hypcap=false,font=small,labelfont=bf]{caption}
\usepackage{subcaption}
\usepackage{tikz}
\usetikzlibrary{arrows.meta,positioning,calc}
\usepackage{pgfplots}
\pgfplotsset{compat=1.18}
\usepackage{enumitem}
\usepackage{xcolor}
\usepackage{array}
\usepackage{siunitx}
\sisetup{detect-all,table-number-alignment=center,
         table-format=1.3,group-digits=false}
\usepackage[framemethod=TikZ]{mdframed}
\usepackage{needspace}
\usepackage{hyperref}
\usepackage[nameinlink,noabbrev]{cleveref}

\graphicspath{{figures/}}

% --- Palette ----------------------------------------------------------------
% One accent, shared with the figures: the #0072B2 of the house matplotlib
% style, so that a coloured rule on the page and a blue curve in a panel are
% the same blue.  \accentpale is the wash used behind highlighted results.
\definecolor{accent}{HTML}{0072B2}
\definecolor{accentpale}{HTML}{EFF6FB}
\definecolor{accentedge}{HTML}{BBD6E9}
\definecolor{ink}{HTML}{1A1C1F}
\definecolor{soft}{HTML}{565B62}

% --- Vertical justification and page breaking -------------------------------
% Never break a paragraph leaving one line stranded at the foot or head of a
% page; prefer a slightly short page to a widow.
\widowpenalty=10000
\clubpenalty=10000
\displaywidowpenalty=10000
\brokenpenalty=10000
\raggedbottom

% --- Equation cross-references as bare parenthesised numbers -----------------
\crefformat{equation}{(#2#1#3)}
\Crefformat{equation}{(#2#1#3)}
\crefrangeformat{equation}{(#3#1#4)--(#5#2#6)}
\crefmultiformat{equation}{(#2#1#3)}{ and (#2#1#3)}{, (#2#1#3)}{ and (#2#1#3)}

% --- Notation ---------------------------------------------------------------
% Scalars are set in italic; bold is reserved for vector-valued states, which
% appear only in the multi-type work of later chapters.  Cell age is
% \alpha throughout: the letter a is a root of the characteristic quadratic
% and nothing else.  The eclipse conversion rate is \omega.
\newcommand{\pgf}{probability generating function}
\newcommand{\rv}{random variable}
\newcommand{\pr}[1]{\mathbb{P}\!\left\{#1\right\}}
\newcommand{\ex}[1]{\mathbb{E}\!\left(#1\right)}
\newcommand{\mean}[1]{\mathbb{E}\!\left[#1\right]}
\newcommand{\var}[1]{\operatorname{Var}\!\left(#1\right)}
\newcommand{\dd}{\mathrm{d}}
\newcommand{\ee}{\mathrm{e}}
\newcommand{\ddt}[1]{\frac{\dd #1}{\dd t}}
\newcommand{\dda}[2]{\frac{\dd #1}{\dd #2}}
\newcommand{\pddt}[1]{\frac{\partial #1}{\partial t}}
\newcommand{\pdd}[2]{\frac{\partial #1}{\partial #2}}
\newcommand{\lrb}[1]{\left(#1\right)}
\newcommand{\lrbs}[1]{\left[#1\right]}
\newcommand{\ind}{\mathbf 1}

% --- Chapter-specific objects -----------------------------------------------
% Intracellular: birth \lambda, death \mu, catastrophe \delta; absorbing
% rupture state H; released count W; burst size \Kb.  Population: \Icell
% infected cells, \Vfree free particles.
\newcommand{\Ifix}[1][]{I_{\mathrm{fix}#1}}
\newcommand{\Kb}{\mathcal{K}}
\newcommand{\Icell}{\mathcal{I}}
\newcommand{\Vfree}{\mathcal{V}}
\newcommand{\Lap}[1]{\widetilde{#1}}
\newcommand{\Fhyp}{{}_2F_1}
\newcommand{\ft}{\textit{F.~tularensis}}
\newcommand{\yp}{\textit{Y.~pestis}}

% --- Cross-chapter references -----------------------------------------------
% Redefining these macros is all that thesis integration requires: they
% become "Chapter~\ref{ch:...}" once the chapter order is fixed.  No chapter
% number appears anywhere in the body prose.
\newcommand{\ChCore}{the chapter on the birth--death--catastrophe process}
\newcommand{\ChCorecap}{The chapter on the birth--death--catastrophe process}
\newcommand{\ChDist}{the chapter on its distribution theory}
\newcommand{\ChDistcap}{The chapter on its distribution theory}
\newcommand{\ChTwoType}{the two-type chapter}
\newcommand{\ChTwoTypecap}{The two-type chapter}
\newcommand{\ChPath}{the chapter on bursting and budding pathogenesis}
\newcommand{\ChEvo}{the chapter on variability in evolution}

% --- Forward references to later modelling chapters -------------------------
% Every forward reference is routed through \fwd, whose first argument is a
% stable key and whose second is the prose actually typeset.
\newcommand{\fwd}[2]{#2}

% --- Theorem environments ---------------------------------------------------
% Two visual weights.  The chapter turns on four results -- the projection,
% the invariance of R_0, the flooding advantage and the coincidence at L=1 --
% and a reader flipping through should be able to find them.  `theorem` is
% therefore set apart with a tinted panel and an accent rule; the supporting
% propositions, definitions and remarks stay in the ordinary run of text so
% that the distinction means something.
\newtheoremstyle{ch6plain}
  {0.9\baselineskip}{0.9\baselineskip}   % space above / below
  {\itshape}{}                            % body font, indent
  {\bfseries\color{accent}}{.}            % head font, punctuation
  {0.5em}{}                               % head/body sep, head spec
\newtheoremstyle{ch6definition}
  {0.9\baselineskip}{0.9\baselineskip}
  {\normalfont}{}
  {\bfseries}{.}
  {0.5em}{}

\theoremstyle{ch6plain}
\newtheorem{theorem}{Theorem}[chapter]
\newtheorem{proposition}[theorem]{Proposition}
\newtheorem{lemma}[theorem]{Lemma}
\newtheorem{corollary}[theorem]{Corollary}
\theoremstyle{ch6definition}
\newtheorem{definition}[theorem]{Definition}
\theoremstyle{remark}
\newtheorem{remark}[theorem]{Remark}

% The tinted panel around `theorem` only.
\mdfdefinestyle{ch6thm}{
  linewidth=0pt,
  leftline=true,linecolor=accent,linewidth=1.6pt,
  topline=false,bottomline=false,rightline=false,
  backgroundcolor=accentpale,
  innerleftmargin=0.85em,innerrightmargin=0.85em,
  innertopmargin=0.7em,innerbottommargin=0.7em,
  skipabove=0.9\baselineskip,skipbelow=0.9\baselineskip,
  nobreak=false,
}
\surroundwithmdframed[style=ch6thm]{theorem}

% Proofs: a shade quieter than the surrounding prose, so a statement stands
% out from its own argument.
\renewenvironment{proof}[1][\proofname]{%
  \par\vspace{0.25\baselineskip}\noindent
  \small\color{soft}%
  \textit{#1.}\ \ignorespaces}{%
  \unskip\nobreak\hfill$\square$\par\vspace{0.35\baselineskip}}

% --- A highlighted standalone result ----------------------------------------
% For the one equation the chapter is built around.  Used once, in the
% flooding section; using it twice would spend the emphasis.
\newmdenv[
  linewidth=0.9pt,linecolor=accentedge,
  backgroundcolor=accentpale,
  roundcorner=2.5pt,
  innerleftmargin=1em,innerrightmargin=1em,
  innertopmargin=0.55em,innerbottommargin=0.55em,
  skipabove=1.0\baselineskip,skipbelow=1.0\baselineskip,
]{keyresult}

\allowdisplaybreaks
